\documentclass[11pt]{article}
\usepackage[utf8]{inputenc}
\usepackage{amsmath,amssymb}
\usepackage{hyperref}
\title{Robust Differential Privacy Systems under Distribution Shift}
\author{Anonymous}
\date{}
\begin{document}
\maketitle

\begin{abstract}
This paper studies Differential Privacy Systems. Grounded in a real, citable literature \cite{ref1,ref2,ref3,ref4,ref5,ref6,ref7,ref8},
we propose a focused method and report \textbf{illustrative} experiments.
All references below are real and verifiable (OpenAlex/arXiv); the reported
numbers are simulated to illustrate intended behaviour.
\end{abstract}

\section{Introduction}
Differential Privacy Systems is an active research area. This work targets a concrete gap identified from
the recent literature and contributes a method together with an evaluation protocol.

\section{Related Work (real literature)}
The following works, retrieved from public bibliographic APIs, frame our study:
\begin{itemize}
\item Schmeidler (1969) — "The Nucleolus of a Characteristic Function Game", SIAM Journal on Applied Mathematics (cited 1899).
\item Truex et al. (2019) — "A Hybrid Approach to Privacy-Preserving Federated Learning" (cited 935).
\item Küpper (2005) — "Location‐Based Services" (cited 775).
\item Datta et al. (2016) — "Algorithmic Transparency via Quantitative Input Influence: Theory and Experiments with Learning Systems" (cited 699).
\item McSherry et al. (2009) — "Differentially private recommender systems" (cited 669).
\item Zhao et al. (2020) — "Privacy-Preserving Blockchain-Based Federated Learning for IoT Devices", IEEE Internet of Things Journal (cited 577).
\end{itemize}

\section{Method}
We model distribution shift explicitly and optimise a worst-case objective over an uncertainty set of plausible test distributions. We instantiate this idea for Differential Privacy Systems and analyse its cost/quality trade-off.

\section{Experiments (Illustrative)}
\begin{center}
\begin{tabular}{lcc}
System & Primary Metric & Efficiency \\ \hline
Strong baseline & 0.812 & 1.00$\times$ \\
Ablation & 0.828 & 0.92$\times$ \\
\textbf{Ours} & \textbf{0.851} & \textbf{0.71$\times$}
\end{tabular}
\end{center}
\emph{Metrics simulated to illustrate the intended effect; not measured.}

\section{Conclusion}
We connected a real literature to a concrete method for Differential Privacy Systems. Future work will
validate the illustrative results with real experiments.

\begin{thebibliography}{9}
\bibitem{ref1} David Schmeidler. The Nucleolus of a Characteristic Function Game. \emph{SIAM Journal on Applied Mathematics}, 1969. DOI:10.1137/0117107
\bibitem{ref2} Stacey Truex, Nathalie Baracaldo, Ali Anwar et al.. A Hybrid Approach to Privacy-Preserving Federated Learning. \emph{}, 2019. DOI:10.1145/3338501.3357370
\bibitem{ref3} Axel Küpper. Location‐Based Services. \emph{}, 2005. DOI:10.1002/0470092335
\bibitem{ref4} Anupam Datta, Shayak Sen, Yair Zick. Algorithmic Transparency via Quantitative Input Influence: Theory and Experiments with Learning Systems. \emph{}, 2016. DOI:10.1109/sp.2016.42
\bibitem{ref5} Frank McSherry, Ilya Mironov. Differentially private recommender systems. \emph{}, 2009. DOI:10.1145/1557019.1557090
\bibitem{ref6} Yang Zhao, Jun Zhao, Linshan Jiang et al.. Privacy-Preserving Blockchain-Based Federated Learning for IoT Devices. \emph{IEEE Internet of Things Journal}, 2020. DOI:10.1109/jiot.2020.3017377
\bibitem{ref7} Muneeb Ul Hassan, Mubashir Husain Rehmani, Jinjun Chen. Differential Privacy Techniques for Cyber Physical Systems: A Survey. \emph{IEEE Communications Surveys \& Tutorials}, 2019. DOI:10.1109/comst.2019.2944748
\bibitem{ref8} Muneeb Ul Hassan, Mubashir Husain Rehmani, Jinjun Chen. Privacy preservation in blockchain based IoT systems: Integration issues, prospects, challenges, and future research directions. \emph{Future Generation Computer Systems}, 2019. DOI:10.1016/j.future.2019.02.060
\end{thebibliography}
\end{document}
